\begin{abstract}
Although large unsupervised language models (LMs) acquire extensive world knowledge and some capacity for reasoning, their behavior remains challenging to direct precisely because their training lacks any explicit supervision. Current strategies for enhancing the controllability of LMs typically involve gathering human preference data by ranking generated responses, then fine-tuning the model—most commonly through reinforcement learning from human feedback (RLHF)—to align output with these preferences. Yet, RLHF is inherently intricate and can be unstable: it first constructs a reward model that approximates human preferences, then further adjusts the language model using reinforcement learning to increase this reward signal while preserving proximity to the initial distribution. This work introduces a novel way to parametrize the reward model within RLHF, making it possible to derive the optimal policy analytically and thus recast the standard RLHF objective as a straightforward classification task. We refer to this approach as \textit{Direct Preference Optimization} (DPO). DPO is robust, achieves strong performance, and significantly reduces computational complexity; it obviates the requirement for in-loop sampling or intensive hyperparameter optimization during fine-tuning. Our empirical results demonstrate that DPO matches or surpasses the effectiveness of established approaches in aligning LMs with human preferences. In particular, DPO achieves greater sentiment control compared to RLHF with PPO, and yields comparable or superior results for tasks such as summarization and single-turn dialogue—while being considerably easier to train and implement.
\end{abstract}

\section{Introduction}
Recent advances have shown that large unsupervised language models (LMs), when trained on vast corpora, can develop unexpectedly sophisticated capabilities~\citep{chowdhery2022palm, brown2020language, touvron2023llama, bubeck2023sparks}. Despite these achievements, such models are built on human-generated data characterized by highly diverse intentions, objectives, and levels of expertise. Not all such characteristics are worth emulating; for instance, an AI-based programming assistant should \textit{recognize} common coding mistakes to remedy them, yet, ideally, its code generation ought to prioritize the (potentially infrequent) examples of superior coding proficiency found within its training data. In a similar vein, a language model should be \textit{informed} of a widespread misconception held by half the population, but it is clearly undesirable for the model to replicate this error in half of its responses to relevant queries. Thus, distinguishing and reinforcing the \emph{preferred behaviors and outputs} from the model’s extensive \textit{repertoire of knowledge and skills} is fundamental to constructing AI systems that are robust, aligned, and manageable~\citep{ouyang2022training}. Although the predominant techniques currently align LMs with human preferences through reinforcement learning (RL), we will demonstrate that the objective commonly pursued by RL-based methods can in fact be solved precisely via a straightforward binary cross-entropy formulation, which streamlines the overall preference alignment workflow.

Broadly, prevailing approaches integrate the target behaviors into LMs through meticulously selected collections of human preference data exemplifying what users deem helpful and trustworthy conduct. This process, often called preference learning, is typically performed following an initial phase of massive-scale, unsupervised training on text data. While preference alignment could simply use supervised fine-tuning on curated high-quality human demonstrations, the most effective strategies to date use reinforcement learning from human (or AI) feedback (RLHF/RLAIF; \citep{christiano2017deep,bai2022constitutional}). RLHF methods involve fitting a reward model on a set of human preferences and then applying RL to optimize the language model’s policy to elicit outputs that receive high reward, all while maintaining sufficient similarity to the original, pre-trained model. Although RLHF produces language models with remarkable ability in conversation and code generation, the overall RLHF methodology introduces considerable complexity compared to supervised learning. It requires the training of several LMs and continual sampling from the policy during the training process, thereby imposing notable computational burdens.

This paper presents a method to directly optimize language models according to human preferences, foregoing explicit reward modeling and reinforcement learning techniques. We introduce \textit{Direct Preference Optimization (DPO)}, an approach that implicitly targets the same objective as traditional RLHF algorithms—namely, maximizing reward subject to a KL-divergence penalty—but does so through a more streamlined and accessible training process. In essence, DPO modifies the log probability ratio between preferred and non-preferred responses, utilizing an adaptive, example-specific importance weighting mechanism. This mechanism serves to counteract the model deterioration that can result from a simplistic application of the probability ratio. Similar to existing frameworks, DPO employs a theoretical model of preferences, such as the Bradley-Terry model \cite{bradley1952rankanalysis}, to quantify the consistency between a reward function and observed human preferences. Whereas conventional methods leverage this model to establish a loss for training a reward model—subsequently used to guide policy learning—DPO reparametrizes the preference loss, allowing it to directly operate as a function of the policy itself. Given a dataset containing pairwise human preferences among responses, DPO can optimize a policy by minimizing a simple binary cross entropy loss, thereby discovering the policy that optimally aligns with a latent reward inferred from the preference data.

The primary contribution of our work is Direct Preference Optimization (DPO), an RL-free and straightforward approach to learning from human preferences in language model training. Empirically, we demonstrate that DPO matches or exceeds the performance of established methods, including PPO-based RLHF, across a range of tasks—such as sentiment control, text summarization, and dialogue generation—using models with up to 6B parameters.

\section{Related Work}

Self-supervised language models with larger parameter counts have demonstrated the capacity to solve certain tasks in the absence of explicit supervision—either in a zero-shot manner \citep{radford2019language} or through few-shot prompting \citep{gpt3,megatron,chowdhery2022palm}. Nevertheless, their effectiveness on downstream tasks and the degree to which their outputs align with user intentions can be substantially enhanced by fine-tuning with datasets comprising instructions and human-authored responses \citep{mishra-etal-2022-cross,sanh2022multitask,chung2022scaling,thoppilan2022lamda}. This technique, known as ‘instruction-tuning’, fosters better generalization of LLMs to follow novel instructions not present in the original instruction-tuning set, thereby improving overall user experience \citep{chung2022scaling}. Although instruction tuning has achieved notable success, collecting \textit{relative} assessments of response quality from human annotators is often more feasible than curating expert-labeled examples. Consequently, later research has relied on human preference datasets to further fine-tune LLMs, resulting in increased skill in translation \citep{kreutzer-etal-2018-reliability}, summarization \citep{stiennon2022learning,ziegler2020finetuning}, narrative generation \citep{ziegler2020finetuning}, and instruction adherence \citep{ouyang2022training,ramamurthy2023is}. In such methods, a neural network reward model is initially trained to align with these preference datasets using a preference structure like the Bradley-Terry model \citep{bradley1952rankanalysis}; the language model is then further optimized to maximize the estimated reward, commonly employing reinforcement learning strategies such as REINFORCE \citep{williams1992reinforce}, proximal policy optimization (PPO; \cite{schulman2017proximal}), or their variants \citep{ramamurthy2023is}. A related research direction exploits LLMs refined via instruction-following with human-in-the-loop feedback to generate supplementary synthetic preference data for specific criteria—such as safety or non-harmfulness \citep{bai2022constitutional}—with weak supervision provided only by a text-based rubric for annotation. This convergence integrates research traditions from training language models with reinforcement learning to meet diverse objectives~\citep{Ranzato2015SequenceLT,paulus2018a,wu2018learning} and from developing general approaches for learning from human feedback \citep{christiano2017deep,kupcsik2018learning}. While the prospect of utilizing relative human preferences is compelling, practical limitations still impede reinforcement learning-based fine-tuning of large language models; the present work introduces a theoretically principled strategy for optimizing with relative preferences in the absence of RL.

Beyond the language domain, the task of learning policies from preferences has been extensively investigated within both bandit and reinforcement learning frameworks, leading to the development of multiple methodologies. When preferences or action rankings, as opposed to explicit rewards, guide contextual bandit learning, the problem is framed as a contextual dueling bandit (CDB; \cite{yue2012karmed,dudik2015contextual}). In scenarios where absolute rewards are not observable, theoretical work on CDBs introduces the \textit{von Neumann winner}: a policy whose probability of outperforming any alternative policy is at least 50\% in expectation \citep{dudik2015contextual}. A key distinction in the CDB literature is that preferences are provided in an online fashion, whereas, for human preference learning, one generally relies on a static, offline dataset comprised of preference-annotated pairs of actions \citep{yan2022human}. Analogously, \textit{preference-based RL} (PbRL) is concerned with learning from binary preferences that arise from an unspecified `scoring' function, in lieu of observed rewards \citep{BusaFekete2014,ruiz2023dueling}. There exists a diversity of algorithms for PbRL, some of which facilitate the reuse of off-policy preference data; these generally require an explicit intermediate step: inferring the hidden scoring function—effectively a reward model—followed by its optimization \citep{jain2013learning,BusaFekete2014,christiano2017deep,sadigh2017active,kupcsik2018learning}. In contrast, our approach dispenses with this separation, proposing an integrated policy learning method that directly adjusts the policy to conform to given preferences.

\section{Preliminaries}\label{section:prelims}

We provide an overview of the RLHF pipeline as originally described by \citeauthor{ziegler2020finetuning} and subsequently refined in works such as \citep{stiennon2022learning, bai2022training, ouyang2022training}. The process is generally partitioned into three principal stages: (1) supervised fine-tuning (SFT); (2) reward modeling based on preference data; and (3) optimization using reinforcement learning.

\textbf{SFT}: The first step in RLHF involves adapting a pre-trained language model (LM) via supervised learning, utilizing curated, high-quality data pertinent to the target application (e.g., dialogue, summarization). This results in a fine-tuned model $\pi^\text{SFT}$.

\textbf{Reward Modelling Phase}: During the subsequent stage, the SFT model is exposed to prompts $x$ to generate answer pairs $(y_1, y_2)\sim \pi^\text{SFT}(y \mid x)$. These answer pairs are assessed by human annotators who indicate a preference for one response, denoted by $y_w\succ y_l \mid x$, where $y_w$ and $y_l$ identify the preferred and non-preferred responses, respectively, from the pair $(y_1, y_2)$. The presumption is that these preferences are governed by some unobserved reward function $r^*(y, x)$, to which we do not have direct access. Multiple models exist to represent preferences, with the Bradley-Terry (BT) model \cite{bradley1952rankanalysis} being a prevalent choice. Alternatively, more general ranking models such as Plackett-Luce \citep{plackett1975analysis, luce2012individual} are also applicable, especially when human rankings are provided for multiple completions. According to the BT model, the distribution over human preferences $p^*$ is formulated as follows:

\begin{equation}\label{eq:bradley-terry}
    p^*(y_1\succ y_2 \mid x)=\frac{\exp\left(r^*(x, y_1)\right)}{\exp\left(r^*(x, y_1)\right) + \exp\left(r^*(x, y_2)\right)}.
\end{equation}

Given a fixed dataset of comparison samples $\mathcal{D}=\bigl\{x^{(i)}, y_w^{(i)}, y_l^{(i)}\bigr\}_{i=1}^N$ drawn from $p^*$, we introduce a parametric reward model $r_{\phi}(x, y)$ and estimate its parameters using the maximum likelihood principle. Treating the setup as a binary classification, the training objective becomes the negative log-likelihood loss:

\begin{equation}\label{eq:reward_model}
    \mathcal{L}_R(r_{\phi}, \mathcal{D}) = -\mathbb{E}_{(x, y_w, y_l)\sim \mathcal{D}}\bigl[\log \sigma(r_{\phi}(x, y_w)- r_{\phi}(x, y_l))\bigr]
\end{equation}

where $\sigma$ denotes the logistic sigmoid. In the language model domain, $r_{\phi}(x, y)$ is commonly initialized using the SFT model $\pi^\text{SFT}(y \mid x)$, and a linear layer is added atop the last transformer layer to output a scalar reward prediction \cite{ziegler2020finetuning}. To reduce variance in the reward function, standard practice normalizes the reward values so that $\mathbb{E}_{x,y\sim \mathcal{D}}\left[r_\phi(x, y)\right] = 0$ for each $x$.

\textbf{RL Fine-Tuning Phase}: When entering the RL phase, the reward function trained previously acts as the evaluative signal for updating the language model. Following earlier methodologies~\citep{jaques2017sequence, jaques2020human}, the policy optimization objective is formalized as

\begin{equation}\label{eq:RL}
\max_{\pi_{\theta}}  \mathbb{E}_{x\sim \mathcal{D}, y\sim \pi_{\theta}(y \mid x)}\bigl[r_{\phi}(x, y)\bigr] - \beta\mathbb{D}_{\textrm{KL}}\bigl[\pi_{\theta}(y\mid x)\mid \mid \pi_\text{ref}(y\mid x)\bigr],
\end{equation}

where the hyperparameter $\beta$ determines how much the policy $\pi_{\theta}$ may differ from the reference distribution $\pi_\text{ref}$, specifically the initial SFT model $\pi^\text{SFT}$. 
Practically, the language model policy $\pi_\theta$ is initialized to match $\pi^\text{SFT}$. This constraint is crucial to prevent the policy from straying beyond the input regions where the reward model is trusted and to preserve diversity in the outputs, avoiding the issue of converging on a narrow set of high-reward responses. Due to the discrete token structure of text generation, this optimization problem cannot be solved directly via backpropagation and therefore relies on reinforcement learning. The prevalent method \citep{ziegler2020finetuning, stiennon2022learning, bai2022training, ouyang2022training} is to set the reward as ${r(x, y) = r_{\phi}(x, y) -\beta (\log \pi_{\theta}(y\mid x) - \log \pi_\text{ref}(y\mid x))}$ and maximize it using Proximal Policy Optimization (PPO) \cite{schulman2017proximal}.

\section{Direct Preference Optimization}\label{sec:DPO}

Recognizing the difficulties in scaling reinforcement learning algorithms to large models for fine-tuning language generators, we aim to develop a straightforward method for optimizing policies based directly on preferences. Unlike traditional RLHF pipelines, which first fit a reward function and subsequently perform RL-based optimization, our approach capitalizes on a reward parameterization that allows for analytic recovery of the corresponding optimal policy, thereby obviating the need for iterative RL updates. As outlined in the following discussion, the core insight lies in leveraging a closed-form mapping from the reward to the optimal policy, which in turn enables transforming loss objectives over reward models into losses defined directly on policies. By using this variable transformation, it becomes unnecessary to explicitly train a separate reward model, yet our optimization still adheres to probabil

\textbf{Derivation of the DPO objective.} Our starting point is the RL objective as presented in prior literature, namely Eq.~\ref{eq:RL}, assuming a generic reward function $r$. As established in previous studies~\citep{peters2007reinforcement, peng2019advantage, korbak2022reinforcement, go2023aligning}, it can be readily demonstrated that the optimal policy maximizing expected rewards under a KL-divergence constraint, as formalized in Eq.~\ref{eq:RL}, is given by:

\begin{equation}\label{eq:op_policy}
    \pi_r(y\mid x) = \frac{1}{Z(x)}\pi_\text{ref}(y\mid x)\exp\left(\frac{1}{\beta}r(x, y)\right),
\end{equation}

Here, the normalization term $Z(x) =\sum_{y}\pi_\text{ref}(y\mid x)\exp\left(\frac{1}{\beta}r(x, y)\right)$ serves as the partition function. Detailed steps can be found in Appendix \ref{app:derivation1}. In practical scenarios, even if one employs an MLE-based reward estimator $r_{\phi}$ for the true reward $r^*$, calculating $Z(x)$ remains computationally costly~\citep{korbak2022reinforcement, go2023aligning}, limiting the utility of this form. To address this, one can manipulate Eq.~\ref{eq:op_policy} to reformulate the reward as a function of the optimal policy $\pi_r$, the reference policy $\pi_\text{ref}$, and the partition function $Z(\cdot)$. Specifically, by applying the logarithm and rearranging terms, we arrive at:

\begin{equation}\label{eq:main_eq}
    r(x,y) =\beta \log \frac{\pi_r(y\mid x)}{\pi_\text{ref}(y\mid x)} + \beta \log Z(x).
\end{equation}

This parameterization can be used for the actual reward $r^*$ and its corresponding optimal policy $\pi^*$. A crucial observation is that the Bradley-Terry model relies exclusively on the reward gap between two candidates: ${p^*(y_1 \succ y_2 \mid x) = \sigma(r^*(x, y_1) - r^*(x, y_2))}$. Replacing $r^*(x,y)$ in the preference model (Eq.~\ref{eq:bradley-terry}) with the parameterization from Eq.~\ref{eq:main_eq}, one sees that the partition function cancels out, allowing us to write the preference probability purely in terms of the optimal policy $\pi^*$ and the reference policy $\pi_\text{ref}$. Consequently, the RLHF-optimal policy $\pi^*$ adheres to the preference modeling relationship:

\begin{equation}\label{eq:objective}
    p^*(y_1\succ y_2 \mid x)=\frac{1}{1 + \exp\left(\beta \log \frac{\pi^*(y_2\mid x)}{\pi_\text{ref}(y_2\mid x)} - \beta \log \frac{\pi^*(y_1\mid x)}{\pi_\text{ref}(y_1\mid x)}\right)}
\end{equation}

A comprehensive derivation can be found in Appendix~\ref{app:derivation2}. Although this statement employs the Bradley-Terry framework, an analogous development holds for general Plackett-Luce models~\citep{plackett1975analysis, luce2012individual}; refer to Appendix~\ref{app:plackett_luce_models} for details.

With a direct expression for the likelihood of human preference data as a function of the optimal policy, rather than the reward, we are now positioned to construct a maximum likelihood training objective for a parameterized policy $\pi_\theta$. This is directly analogous to reward model fitting (see Eq.~\ref{eq:reward_model}), and leads to the policy learning objective:

\begin{equation}\label{eq:optimum_model}
    \mathcal{L}_\text{DPO}(\pi_{\theta}; \pi_\text{ref}) = -\mathbb{E}_{(x, y_w, y_l)\sim \mathcal{D}}\left[\log \sigma \left(\beta \log \frac{\pi_{\theta}(y_w\mid x)}{\pi_\text{ref}(y_w\mid x)} - \beta \log \frac{\pi_{\theta

In this formulation, we learn an implicit reward through a different parameterization, such that its associated optimal policy corresponds directly to $\pi_\theta$. Additionally, our approach mirrors the structure of a reparameterized Bradley-Terry model, thereby inheriting various theoretical guarantees, including consistency provided the preference data distribution satisfies certain criteria \cite{bong2022generalized}. We elaborate further on the theoretical implications of DPO in comparison to related methodologies in Section~\ref{sec:theory}.

\textbf{Mechanism of the DPO update.} To elucidate the operational mechanics of DPO, it is instructive to investigate the gradient of its objective function, $\mathcal{L}_\text{DPO}$. The gradient with respect to $\theta$ is expressed as:

\begin{multline*}\label{eq:gradient}
    \nabla_\theta \mathcal{L}_\text{DPO}(\pi_\theta;\pi_\text{ref}) = \\ -\beta\mathbb{E}_{(x, y_w, y_l) \sim \mathcal{D}} \bigg[\underbrace{\sigma(\hat{r}_\theta(x, y_l) - \hat{r}_\theta (x, y_w))}_\text{greater emphasis when reward ordering is inaccurate}\bigg[\underbrace{\nabla_\theta\log \pi(y_w \mid x)}_\text{favor $y_w$ by boosting likelihood} - \underbrace{\nabla_\theta\log\pi(y_l \mid x)}_\text{discourage $y_l$ by reducing likelihood}\bigg]\bigg],
\end{multline*}

with $\hat{r}_\theta(x, y) = \beta \log \frac{\pi_\theta(y \mid x)}{\pi_\text{ref}(y \mid x)}$ representing the reward implicitly determined by the language model $\pi_\theta$ relative to the reference $\pi_\text{ref}$ (see also Section~\ref{sec:theory} for details). Conceptually, this loss gradient amplifies the probability assigned to the preferred continuations $y_w$, while suppressing that of the non-preferred continuations $y_l$. Crucially, the influence of each sample is modulated according to the extent to which the implicit reward $\hat{r}_\theta$ erroneously favors the less preferred completion, scaled by $\beta$—reflecting the degree of misranking as well as the KL constraint's strictness. Our empirical results indicate that incorporating this weighting scheme is critical; omitting it in a simpler variant results in model degeneration (see Appendix Table~\ref{tab:unlikelihood_generations}).

\textbf{DPO overview.} The standard DPO workflow proceeds as follows: 1) For each prompt $x$, generate completion samples $y_1, y_2$ using $\pi_\text{ref}(\cdot \mid x)$, annotate these with human preference signals to create the preference dataset $\mathcal{D} = \{x^{(i)}, y_w^{(i)}, y_l^{(i)}\}_{i=1}^N$; and 2) train the language model $\pi_\theta$ to minimize the DPO loss, $\mathcal{L}_\text{DPO}$, conditioned on the chosen $\pi_\text{ref}$, the constructed dataset $\mathcal{D}$, and target value $\beta$. 
In applied settings, it is advantageous to leverage existing public preference datasets, instead of re-generating data and collecting human preference annotations from scratch. When such datasets have been sampled with $\pi^\text{SFT}$, we set $\pi_\text{ref} = \pi^\text{SFT}$ by default if it is accessible. In cases where $\pi^\text{SFT}$ is unavailable, $\pi_\text{ref}$ is initialized by maximizing the likelihood of preferred responses ${(x, y_w)}$: specifically, ${\pi_\text{ref} = \argmax_{\pi}\mathbb{E}_{x, y_w \sim \mathcal{D}}\left[\log \pi(y_w \mid x)\right]}$. This approach addresses the potential distributional gap between the (unobserved) true reference distribution and the proxy $\pi_\text{ref}$ used by DPO. Further information regarding implementation specifics and hyperparameter configurations can be found in Appendix~\ref{app:implementation}.

\section{Theoretical Analysis of DPO} This section offers a deeper analysis of the DPO technique, supplies theoretical support, and connects DPO's strengths with limitations commonly observed in actor-critic algorithms employed in RLHF frameworks, such as PPO~\cite{schulman2017proximal}.

\label{sec:theory}

\subsection{Your Language Model Is Secretly a Reward Model} DPO eliminates the need for both explicit reward model fitting and reinforcement learning procedures, instead utilizing a single maximum likelihood training objective. Importantly, the optimization goal in Eq. \ref{eq:main_eq} is equivalent to that of a Bradley-Terry model, with rewards parameterized as $r^*(x, y) = \beta \log\frac{\pi^*_\theta(y \mid x)}{\pi_\text{ref}(y \mid x)}$. Here, the parametric model $\pi_{\theta}$ is updated in a manner comparable to optimizing a reward model as shown in Eq. \ref{eq:reward_model}, after performing a suitable change of variables. In the subsequent discussion, we formalize the underlying theory motivating this parameterization, demonstrate that this technique does not restrict the expressivity of reward models that can be represented, and argue that it permits exact recovery of the optimal policy. Our exposition starts by introducing an equivalence relation between reward functions.

\begin{definition}
Two reward functions $r(x, y)$ and $r'(x, y)$ are said to be equivalent if and only if there exists a function $f$ such that ${r(x, y)-r'(x, y) = f(x)}$.
\end{definition}

This definition clearly establishes an equivalence relation, thereby dividing the collection of reward functions into equivalence classes. The following two lemmas can then be formalized:

\begin{lemma}\label{lemma:same_prefrence} Within the Plackett-Luce framework—and, specifically, the Bradley-Terry model—reward functions that belong to the same equivalence class produce identical preference distributions.
\end{lemma}

\begin{lemma}\label{lemma:same_policy}
    For the constrained RL setting, reward functions that are members of the same equivalence class will result in the same optimal policy.
\end{lemma}

Since the proofs are routine, we provide them in Appendix \ref{app:lemma1}. The initial lemma addresses a classical issue of non-identifiability in Plackett-Luce-type models \cite{plackett1975analysis}. This ambiguity necessitates the introduction of extra identifiability constraints to ensure well-defined MLE solutions for Eq.~\ref{eq:reward_model} \cite{bong2022generalized}. The second lemma guarantees that any reward function from a given equivalence class leads to an identical optimal policy. Thus, in constructing our final objective, it suffices to recover any single representative from the target equivalence class. We establish the following theorem, whose proof appears in Appendix~\ref{app:thm1}:

\begin{theorem}\label{thm:main}
    Assuming mild conditions, every reward class compatible with the Plackett-Luce (including Bradley-Terry) models admits a representation of the form ${r(x, y) = \beta \log \frac{\pi(y\mid x)}{\pi_\text{ref}(y\mid x)}}$ for a suitable model $\pi(y\mid x)$ and some specified reference model $\pi_\text{ref}(y \mid x)$.
\end{theorem}

\begin{sproof}
    Let $r(x, y)$ be an arbitrary reward function that induces the optimal model $\pi_r(y \mid x)$ as per Eq.~\ref{eq:op_policy}. We aim to show that some element of the equivalence class associated with $r$ can be cast using the reparameterization above. To do so, define the projection operator $f$ as
\begin{equation}
    f(r; \pi_\text{ref}, \beta)(x, y) = r(x, y) - \beta\log\sum_{y}\pi_\text{ref}(y\mid x)\exp\left(\frac{1}{\beta}r(x, y)\right)
\end{equation}
Here, $f$ effectively adjusts the reward by the log partition function corresponding to $\pi_r$. Since the normalization depends only on $x$, $f(r; \pi_\text{ref}, \beta)(x, y)$ remains within the same equivalence class as $r(x, y)$. Substituting the right-hand side of Eq.~\ref{eq:main_eq} (valid for any reward function), we obtain $f(r; \pi_\text{ref}, \beta)(x, y) = \beta \log \frac{\pi_r(y\mid x)}{\pi_\text{ref}(y\mid x)}$. Thus, $f$ constructs a representative from $r$'s equivalence class having the desired reparameterized form, thereby preserving full generality in our modeling of reward functions.
\end{sproof}

An alternative interpretation of Theorem~\ref{thm:main} is that it uniquely determines, within each reward equivalence class, the specific reward function selected by the DPO reparameterization—namely, the one that fulfills

\begin{equation}\label{eq:lag_p}
     \sum_{y}\underbrace{\pi_\text{ref}(y\mid x)\exp\left(\frac{1}{\beta}r(x, y)\right)}_{=\pi(y\mid x)\text{, as given by Thm.~\ref{thm:main} reparam.}} = 1,
\end{equation}

ensuring that $\pi(y\mid x)$ is indeed a properly normalized probability distribution (i.e., all probabilities are nonnegative and sum to unity). Furthermore, from the perspective of Eq.~\ref{eq:op_policy}, it becomes evident that Eq.~\ref{eq:lag_p} corresponds to the partition function for the optimal policy determined by the reward function $r(x, y)$. The central insight behind the DPO methodology is the imposition of specific constraints on the Plackett-Luce family of preference models (with a focus on Bradley-Terry models), which are otherwise underdetermined. These constraints maintain the diversity of permissible reward representations while ensuring that the optimal policy described by Eq.~\ref{eq:op_policy} remains analytically solvable for any given prompt $x$.

\subsection{Instability of Actor-Critic Algorithms}
Our framework further provides a lens through which to analyze the sources of instability commonly observed in standard actor-critic methods employed in RLHF pipelines, such as PPO. Restricting our attention to the RL fine-tuning stage, as detailed in Section \ref{section:prelims}, we draw parallels with the control as inference formulation presented by \cite{levine2018reinforcement}, focusing on the constrained reinforcement learning problem given in \ref{eq:RL}. Specifically, with a parameterized distribution $\pi_{\theta}(y\mid x)$, optimization proceeds by minimizing the Kullback-Leibler divergence $\mathbb{D}_{\text{KL}}[\pi_{\theta}(y|x) \mid\mid \pi^*(y\mid x)]$, where $\pi^*$ represents the reward-induced optimal policy of Eq.~\ref{eq:optimum_model} for reward function $r_{\phi}(y, x)$. Elementary manipulations yield the following objective for optimization:

\begin{equation}\label{eq:AC}
    \max_{\pi_{\theta}}\mathbb{E}_{\pi_{\theta}(y\mid x)}\left[\underbrace{r_{\phi}(x, y) -\beta\log\sum_{y}\pi_\text{ref}(y\mid x)\exp\left(\frac{1}{\beta}r_{\phi}(x, y)\right)}_{f(r_{\phi}, \pi_\text{ref}, \beta)} - \underbrace{\beta\log\frac{\pi_{\theta}(y\mid x)}{\pi_\text{ref}(y\mid x)}}_{\text{KL}}\right]
\end{equation}

This objective has also been the focus of optimization in previous research \citep{ziegler2020finetuning, stiennon2022learning, bai2022training, ouyang2022training}, employing the DPO-equivalent reward formulation for the reward class $r_{\phi}$. Within this framework, the normalization component in $f(r_{\phi}, \pi_\text{ref}, \beta)$ can be viewed as representing the soft value function of the reference policy $\pi_\text{ref}$. Although the presence or absence of this normalization does not influence the location of the optimum, omitting it can cause the policy gradient estimator for the objective to display significant variance, thereby reducing the stability of learning. A learned value function can be used to handle the normalization term, but optimizing such a function poses its own set of challenges. Another common approach in earlier work has been to normalize the rewards using a baseline derived from human completions, effectively yielding a Monte-Carlo estimate from a single sample for the normalization. By contrast, the DPO reparameterization leads to a reward function that eliminates the need for baseline normalization entirely.

\section{Experiments}
Here, we conduct empirical investigations of DPO’s effectiveness in learning policies directly from preference data. To begin, in a controlled text-generation environment, we examine DPO’s sample efficiency in balancing the objectives of reward maximization and KL-divergence minimization with respect to the reference policy, benchmarking against standard preference optimization algorithms such as PPO. Subsequently, we assess DPO across larger-scale models and more challenging RLHF scenarios, including summarization and dialogue generation tasks. Our results demonstrate that DPO typically matches or outperforms robust baselines, such as RLHF utilizing PPO or selecting the best out of $N$ sampled trajectories according to a learned reward model, and this holds true even with minimal hyperparameter adjustment. Prior to presenting these findings, we outline the experimental procedures; further information is provided in Appendix~\ref{app:exp_details}.

\textbf{Tasks.} Our study investigates three distinct open-ended text generation tasks. In every experiment, models derive a policy using a preference dataset $\mathcal{D}=\bigl\{x^{(i)}, y_w^{(i)}, y_l^{(i)}\bigr\}_{i=1}^N$. For \textbf{controlled sentiment generation}, $x$ corresponds to the beginning of a movie review sourced from the IMDb dataset \cite{maas-EtAl:2011:ACL-HLT2011}, and the model is tasked with producing a continuation $y$ characterized by positive sentiment. To facilitate systematic evaluation, for this task we \textit{construct} pairs of preferred generations automatically, relying on outputs from a pre-trained sentiment classifier, such that $p(\text{positive}\mid x,y_w)>p(\text{positive}\mid x,y_l)$. For supervised fine-tuning (SFT), we adapt GPT-2-large via fine-tuning on the training partition of IMDb movie reviews, continuing until convergence (additional information can be found in App~\ref{app:sentiment_details}). In the \textbf{summarization} task, $x$ is a Reddit forum post, and the model must produce a summary $y$ highlighting the core points. Consistent with previous research, we use the Reddit TL;DR summarization dataset \citep{volske-etal-2017-tl}, accompanied by human preference data gathered in \citeauthor{stiennon2022learning}. Here, the SFT model is further adapted using human-authored summaries of forum posts,\footnote{\url{https://huggingface.co/CarperAI/openai_summarize_tldr_sft}} implemented within the TRLX \citep{leandro_von_werra_2023_7790115} framework for RLHF. Human preference judgments were sourced by \citeauthor{stiennon2022learning} from samples of another SFT model that, while different, underwent similar training. The third setting, \textbf{single-turn dialogue}, presents $x$ as a human inquiry—ranging from scientific questions to personal advice—where the model must supply an informative and engaging response $y$. For this, we utilize the Anthropic Helpful and Harmless dialogue dataset \citep{bai2022training}, which consists of 170k conversational exchanges between a person and an automated assistant, with each ending in two possible model-generated responses and a human preference label indicating which is favored. Since a specialized SFT model is not pre-existing for this context, we construct one by fine-tuning a readily available language model exclusively on the preferred responses from the dataset.

\textbf{Evaluation.} We employ two distinct strategies for evaluating our methods. For the controlled sentiment generation task, where the true reward function (provided by a sentiment classifier) is accessible, we assess each algorithm by mapping out the trade-off curve between achieved reward and KL-divergence relative to the reference policy; this trade-off is feasible to compute in this context. Conversely, in practical scenarios where the ground-truth reward function is unavailable, we resort to measuring the \textit{win rate} of each method compared to a baseline policy. In these cases—namely, summarization and single-turn dialogue—we rely on GPT-4 as a stand-in for human assessment to judge summary quality and the helpfulness of responses. For the summarization domain, the baseline comprises the test set’s reference summaries; in the dialogue setting, the baseline is the human-preferred response from the test set. Despite prior research indicating that large language models may surpass conventional metrics as automated evaluators \citep{Chen2023ExploringTU}, we further support our decision to use GPT-4 through a dedicated human study, detailed in Sec.~\ref{sec:human-judgments}. Our findings demonstrate a strong alignment between GPT-4's evaluations and human judgments, with human-GPT-4 agreement matching or exceeding agreement between independent human annotators.

\textbf{Methods.} Alongside DPO, we benchmark a range of existing techniques aimed at aligning language model behavior with human preferences. As straightforward baselines, we utilize zero-shot prompting of \textbf{GPT-J} \citep{gpt-j} for summarization and two-shot prompting of \textbf{Pythia-2.8B} \citep{biderman2023pythia} for dialogue. We also examine an \textbf{SFT} model, as well as a \textbf{Preferred-FT} variant, which is obtained via supervised fine-tuning on the favored output $y_w$—drawn from the SFT model (for sentiment control and summarization) or a general language model (for dialogue). Another approach under consideration is \textbf{Unlikelihood}~\citep{welleck2019neural}, which trains the model to simultaneously increase the likelihood of $y_w$ and decrease the likelihood of $y_l$, optionally weighted by a tunable coefficient $\alpha \in [0,1]$ on the negative term. Furthermore, we test \textbf{PPO} \citep{schulman2017proximal}, employing a reward model inferred from preference data, as well as \textbf{PPO-GT}, which has access to the authentic reward function and thus serves as an oracle in sentiment experiments. Two variants of PPO-GT are utilized in the sentiment experiments: a standard off-the-shelf implementation \cite{leandro_von_werra_2023_7790115} and a custom-tuned version with reward normalization and hyperparameter optimization, with similar tuning applied for PPO when using learned rewards. Lastly, we incorporate the \textbf{Best of $N$} strategy, in which $N$ responses are sampled from the SFT (or from Preferred-FT in the dialogue task), and the top response—according to a reward model trained on preference annotations—is selected. While this decouples the reward model’s reliability from PPO optimization and often yields strong results, its practicality is hindered by the computational cost of sampling $N$ completions per prompt during evaluation.

\subsection{How well can DPO optimize the RLHF objective?}

The reward maximization objective employed in standard RLHF approaches, which incorporates a KL constraint, is designed to promote reward-seeking behavior while simultaneously limiting divergence from a baseline reference policy. As a result, meaningful evaluation of such methods necessitates consideration of both obtained rewards and the corresponding KL divergence; a method attaining marginally higher reward but with greatly increased KL is not necessarily preferable. In Figure~\ref{fig:frontier-tldr-main}, we plot the trade-off between reward and KL (reward-KL frontier) for several algorithms under the sentiment task. For each method, we conduct a series of training runs, each time varying a hyperparameter governing policy conservativeness: PPO uses a target KL $\in\{3,6,9,12\}$, unlikelihood methods utilize $\beta \in \{0.05,0.1,1,5\}$, and for another variant, $\alpha\in\{0.05,0.1,0.5,1\}$. For preferred-FT, different random seeds are employed. This parameter sweep comprises 22 training runs altogether. Every 100 training steps until models converge, policies are evaluated on a fixed set of test prompts to calculate both their mean reward under the ground-truth reward function and their mean sequence-level KL\footnote{Here, the sum of per-timestep KL-divergences is taken.} relative to the reference policy, $\text{KL}\left(\pi\mid \mid \pi_\text{ref}\right)$. Our results reveal that DPO achieves a frontier that is substantially more favorable than the alternatives, attaining the highest reward while maintaining low KL. This finding stands out for several reasons. Notably, while both DPO and PPO optimize the same underlying objective, DPO demonstrates markedly greater efficiency, as its reward/KL trade-off universally surpasses that of PPO. Furthermore, even in scenarios where PPO is provided with access to true reward values (PPO-GT), DPO is able to achieve a superior reward-KL frontier.

\subsection{Can DPO scale to real preference datasets?}
\label{sec:dpo-real-datasets}
We proceed to assess DPO's fine-tuning capabilities in the contexts of summarization and single-turn dialogue tasks. In the domain of summarization, automatic metrics like ROUGE are known to exhibit weak alignment with human judgments~\citep{stiennon2022learning}, motivating previous research to leverage PPO-based fine-tuning guided by human preferences, which has shown improvements in summary quality. Our experimental protocol involves generating model outputs on the test split of the TL;DR summarization dataset, and then calculating the mean win rate when these outputs are compared against reference summaries from the test data. For all methods evaluated, outputs are sampled at temperatures between 0.0 and 1.0, with the corresponding win rates presented in Figure~\ref{fig:frontier-tldr-main} (right panel). DPO, PPO, and Preferred-FT each perform fine-tuning on the same underlying GPT-J SFT checkpoint\footnote{\url{https://huggingface.co/CarperAI/openai_summarize_tldr_sft}}. Our results show that DPO achieves a win rate near 61\% when sampling at temperature 0.0, surpassing PPO’s best win rate of approximately 57\% at the same temperature. Additionally, DPO secures a higher peak win rate than the optimal $N$-best baseline. It is worth mentioning that hyperparameter tuning for DPO’s $\beta$ was not extensively performed, suggesting these outcomes could be conservative estimates of DPO’s upper bound. Furthermore, DPO demonstrates greater stability across different sampling temperatures compared to PPO, whose effectiveness can fall to that of the original GPT-J model as temperature increases. In contrast, Preferred-FT yields only marginal improvement over the SFT baseline. A direct comparison of DPO and PPO with human raters, discussed in Section~\ref{sec:human-judgments}, reveals that at temperature 0.25, DPO outputs were favored 58\% of the time over PPO samples produced at the same temperature.

For single-turn dialogue tasks, we assess each approach on a portion of the test split from the Anthropic HH dataset \citep{bai2022training} that includes one round of human-assistant exchange. The evaluation with GPT-4 leverages the preferred test completions as references to calculate win rates across the methods. In the absence of a conventional SFT baseline for this context, our process begins with a pre-trained Pythia-2.8B, followed by applying Preferred-FT to train a reference model using selected completions so that the generated completions reflect the same distribution, after which DPO training is conducted. We benchmark these against both the best outcome out of 128 Preferred-FT completions (noting that the Best of $N$ approach reaches its ceiling at 128 samples for this problem, as illustrated in Appendix Figure~\ref{fig:best-of-n}) and a 2-shot prompting setup for the base Pythia-2.8B model, observing that DPO matches or surpasses these alternatives under their optimal temperature settings. Additionally, we examine an RLHF model utilizing PPO, trained on the Anthropic HH dataset \footnote{\url{https://huggingface.co/reciprocate/ppo_hh_pythia-6B}} as provided by a recognized implementation \footnote{\url{https://github.com/CarperAI/trlx/tree/main/examples/hh}}, but are unable to identify any prompting strategy or temperature that outperforms the initial Pythia-2.8B model. Drawing on findings from TL;DR experiments and the observation that Best of 128 and PPO pursue the same reward objective, we treat Best of 128 as an approximate measure of PPO performance. Collectively, our results show that DPO uniquely offers computationally efficient improvements over the Anthropic HH dataset's preferred completions, attaining performance on par with or superior to the more computationally intensive Best of 128 reference. Moreover, Figure~\ref{fig:dialogue-main} demonstrates that DPO reaches optimal performance with relatively few training iterations.

\subsection{Generalization to a new input distribution}

To more thoroughly assess how PPO and DPO perform under distributional shifts, we test the policies learned from the Reddit TL;DR summarization task on a distinct data domain: the test partition of the CNN/DailyMail news article dataset \citep{nallapati-etal-2016-abstractive}. For this evaluation, we use the optimal sampling temperatures previously identified on TL;DR (0 and 0.25). Results can be found in Table~\ref{tab:ood}. To measure summarization quality, we use GPT-4’s win rate relative to dataset gold-standard summaries, adapting the GPT-4 (C) prompt originally used for Reddit TL;DR to reference “news article” instead of “forum post.” On this out-of-distribution data, DPO maintains a clear performance advantage over PPO. These findings serve as preliminary evidence that DPO-trained policies are capable of generalizing to new input distributions comparably to those trained with PPO, despite DPO’s lack of access to extra unlabeled Reddit TL;DR prompts available to PPO.

\subsection{Validating GPT-4 judgments with human judgments}
\label{sec:human-judgments}
A user study was conducted to assess the validity of GPT-4’s preference judgments, drawing on outputs from the TL;DR summarization experiments and incorporating two distinct GPT-4 prompt formulations. The \textbf{GPT-4 (S)} (simple) prompt asks which summary best conveys the key information in the post, whereas the \textbf{GPT-4 (C)} (concise) prompt also incorporates conciseness as a criterion; the latter is evaluated since GPT-4, when using the (S) prompt, tends to favor longer, more redundant summaries than human raters. Full prompt wording can be found in Appendix~\ref{app:prompts}. We perform three pairwise comparisons: the highest-performing method (DPO at temp. 0.25), the lowest-performing (PPO at temp. 1.0), and a mid-range performer (SFT at temp. 0.25), each evaluated against PPO at its best greedy temperature. This sampling ensures coverage of diverse output qualities. With both prompts, we observe that GPT-4’s preferences align with human judgments at rates similar to the inter-human agreement rate, indicating that GPT-4 serves as a reasonable surrogate for human assessment (note: multiple human annotations were obtained for DPO and PPO-1 comparisons only due to resource constraints). The \textbf{GPT-4 (C)} prompt is found to yield win rates more reflective of human preferences and is thus employed for main results in Section~\ref{sec:dpo-real-datasets}. Additional methodological details, including a description of the web interface and the full roster of human evaluators, are provided in Appendix~\ref{app:human-study}.

\section{Discussion}
Preference-based learning presents an effective and scalable approach for developing capable and aligned language models. In this work, we have proposed DPO—a straightforward methodology that enables language model training from preference data without relying on reinforcement learning. DPO circumvents the need to fit preference learning within traditional RL frameworks or leverage conventional RL techniques, instead establishing a direct correspondence between language model policies and reward functions. This connection allows models to be optimized according to human preferences using a basic cross-entropy loss, foregoing reinforcement learning without any sacrifice in generality. Despite requiring minimal hyperparameter adjustment, DPO achieves performance on par with, or superior to, current RLHF methods such as those utilizing PPO. As a result, DPO offers a practical means of reducing obstacles associated with preference-driven language model training.

\textbf{Limitations \& Future Work.} Our findings bring to light several open questions for future exploration. For instance, to what extent does a DPO-trained policy exhibit out-of-distribution generalization relative to models optimized with explicit reward functions? Preliminary experiments indicate that DPO-trained models may generalize in a manner comparable to those trained via PPO, yet a more detailed investigation is warranted. Another avenue is to assess whether using the DPO policy for self-labeling can enable effective exploitation of unlabeled data. Furthermore, understanding the dynamics of reward over-optimization within the DPO framework is necessary, particularly in relation to the minor performance drop observed in Figure~\ref{fig:dialogue-main}-right—could this be a symptom of such over-optimization? Moreover, while our current studies involve models of up to 6B parameters, scaling DPO to substantially larger, cutting-edge models presents a compelling direction for upcoming research. With respect to evaluation, our analyses show that GPT-4-based win rates are sensitive to prompt phrasing; therefore, further work is required to determine optimal methods for obtaining high-quality assessments from automated evaluators. Lastly, DPO has the potential to be applied far beyond language models, opening opportunities for preference-based training in diverse generative modeling domains.

\end{document}